\section{Chapter Preface}
The book closes by returning to its opening question: why does incomplete observation force structured inference? The journey from projection to repair can now be read as one cumulative answer.

Each chapter introduced a new layer—constraints, geometry, propagation, optimization, dynamics, and reconfiguration—needed to transform ambiguity into coherent reconstruction.

The conclusion makes that circular narrative explicit: the final synthesis theorem illuminates the first proposition and anchors CLIO as a unified theory rather than a metaphor set.

\section{Derivations and Formal Results}
\begin{theorem}[CLIO synthesis theorem]
Let \(\mathfrak C=(\stateSpace,\observationSpace,\projection,\constraintSet,F,J)\) be a CLIO system. Assume:
\begin{enumerate}
  \item \(A_y\) is nonempty,
  \item \(A_y\) is invariant under \(F\),
  \item There exists a Lyapunov function strictly decreasing outside \(A_y\).
\end{enumerate}
Then \(x_{t+1}=F(x_t)\) converges toward admissible reconstruction. If \(J\) is continuous and \(A_y\) compact, preferred reconstructions exist. If \(J\) is strictly convex on convex \(A_y\), that preferred reconstruction is unique.
\end{theorem}

\[
\boxed{\text{observation}\to\text{constraint}\to\text{admissibility}\to\text{reconstruction}\to\text{optimization}\to\text{repair when reconstruction fails}}
\]
